% Generated draft: requires mathematical and source review.
\documentclass[11pt]{article}
\usepackage[margin=25mm]{geometry}
\usepackage{fontspec}
\setmainfont{Latin Modern Roman}
\usepackage{amsmath,amssymb,mathtools,unicode-math}
\setmathfont{Latin Modern Math}
\usepackage{enumitem,needspace,xurl,hyperref}
\setlength{\parindent}{0pt}
\setlength{\parskip}{5pt}
\setlength{\emergencystretch}{3em}
\newcommand{\N}{\mathbb{N}}
\newcommand{\Z}{\mathbb{Z}}
\newcommand{\Q}{\mathbb{Q}}
\newcommand{\R}{\mathbb{R}}
\newcommand{\C}{\mathbb{C}}
\newcommand{\F}{\mathbb{F}}
\newcommand{\A}{\mathbb{A}}
\newcommand{\PP}{\mathbb{P}}
\newcommand{\E}{\mathbb{E}}
\newcommand{\eps}{\varepsilon}
\newcommand{\dd}{\,\mathrm{d}}
\begin{document}
\section*{UnsolvedMath ID: 11100058}
\textbf{Source problem code:} catalogue entry
\subsection{Definitions and mathematical statement}
For a model category $\mathcal C$ and a small indexing category $I$, consider the relative diagram category $\mathcal C^I$ with objectwise weak equivalences, and set
\[
 \mathbb D_{\mathcal C}(I)=\operatorname{Ho}(\mathcal C^I).
\]
For $u:I\to J$, restriction induces $u^*:\mathbb D_{\mathcal C}(J)\to\mathbb D_{\mathcal C}(I)$. With the appropriate existence hypotheses, homotopy left and right Kan extensions give adjoints
\[
 u_!\dashv u^*\dashv u_*.
\]
For the projection $I\to *$, these encode homotopy colimits and limits. The system, rather than only $\mathbb D_{\mathcal C}(*)=\operatorname{Ho}(\mathcal C)$, retains homotopy-diagram information. In general $\operatorname{Ho}(\mathcal C^I)$ is not simply the ordinary diagram category $\operatorname{Ho}(\mathcal C)^I$.

Hovey's Model categories item 12 asks whether this enlarged viewpoint, proposed in his book appendix, has value. The task is programmatic: develop and assess the collection of all such homotopy diagram categories with their induced adjunctions so that homotopy limits and colimits belong to the structure. Such systems are commonly organized as homotopy prederivators, with additional axioms giving derivators; the excerpt does not prescribe those axioms or a theorem of reconstruction/equivalence.

The registry relates this viewpoint to modern $\infty$-categories and cites Lurie's \emph{Higher Topos Theory}. This is relevant context, not a definition identifying all the above data with an arbitrary $\infty$-category. Existence of every small limit and colimit is an additional completeness/cocompleteness property, not the definition of an $\infty$-category. The comparison and the precise ambient/size and Kan-extension hypotheses require specification.
\subsection{Short English statement}
Study the system of homotopy categories of all diagram shapes, with restriction and homotopy Kan-extension adjunctions, as richer information than one homotopy category. The proposed viewpoint is closely related to derivator and higher-category structures.
\subsection{Sources}
\begin{itemize}
\item Local registry: \url{data/unsolvedmath/problems.json}, permanent UnsolvedMath ID $11100058$, exact source code \texttt{catalogue entry}, attributed to Mark Hovey.
\item \emph{Hovey's Algebraic Topology Problem List}, Model categories item 12: \url{https://www-users.cse.umn.edu/~tlawson/hovey/}. Supplied access date: 2026-07-29; extraction: html-ordered-list; 65 top-level items in nine category pages.
\item The source refers to the appendix of Hovey's book on model categories without further bibliographic data. The registry cites Lurie, \emph{Higher Topos Theory}, without a precise theorem or location, and its final assessment is truncated.
\end{itemize}
\end{document}
